% ============================================================
% Выходные сведения (оборот титула)
% ============================================================
\pagestyle{plain}
\pagecolor{white}

% Блок с копирайтом, редакцией и датой
\noindent
$\copyright$ \the\year{} Юнител Инжиниринг
\hfill
\parbox{0.3\textwidth}{\raggedleft Редакция: \textcolor{unigray}{\rev}}

\vspace{3mm}
\noindent
Москва
\hfill
\parbox{0.3\textwidth}{\raggedleft Дата: \textcolor{unigray}{\revdate}}
\vspace{8mm}

% --- Определение \device по коду ---
% Если ни одно условие не сработает, оставим пустым (или можно задать значение по умолчанию)
\providecommand{\device}{}

\IfSubStr{ЮТКБ.656122.609 РЭ6}{\devicecode}{%
  \def\device{микропроцессорное устройство защиты и автоматики силового трансформатора напряжением до 35 кВ типа «ЮНИТ-М300-ДЗТ2»}%
}{}%
\IfSubStr{ЮТКБ.656122.609 РЭ8}{\devicecode}{%
  \def\device{микропроцессорное устройство защиты и автоматики силового трансформатора напряжением до 35 кВ типа «ЮНИТ-М300-Т»}%
}{}%
\IfSubStr{ЮТКБ.656122.609 РЭ7}{\devicecode}{%
  \def\device{микропроцессорное устройство защиты и автоматики силового трансформатора напряжением до 35 кВ типа «ЮНИТ-М300-Т2»}%
}{}%
\IfSubStr{ЮТКБ.656122.609 РЭ2}{\devicecode}{%
  \def\device{микропроцессорное устройство защиты и автоматики отходящей линии напряжением до 35 кВ типа «ЮНИТ-М300-ОЛ»}%
}{}%
\IfSubStr{ЮТКБ.656122.609 РЭ5}{\devicecode}{%
  \def\device{микропроцессорное устройство защиты и автоматики трансформатора напряжения 6-35 кВ типа «ЮНИТ-М300-ТН»}%
}{}%
\IfSubStr{ЮТКБ.656122.609 РЭ3}{\devicecode}{%
  \def\device{микропроцессорное устройство защиты и автоматики вводного выключателя 6-35 кВ типа «ЮНИТ-М300-ВВ»}%
}{}%
\IfSubStr{ЮТКБ.656122.609 РЭ4}{\devicecode}{%
  \def\device{микропроцессорное устройство защиты и автоматики секционного выключателя 6-35 кВ типа «ЮНИТ-М300-СВ»}%
}{}%
\IfSubStr{ЮТКБ.656122.609 РЭ9}{\devicecode}{%
  \def\device{микропроцессорное устройство автоматического регулирования напряжения силового трансформатора типа «ЮНИТ-М300-АРНТ»}%
}{}%

% Если \device всё ещё пуст – предупредим в логе (но не останавливаем компиляцию)
\ifx\device\empty
  \PackageWarning{myconfig}{Неизвестный devicecode: \devicecode}%
\fi


\par
Настоящее руководство по эксплуатации (РЭ) распространяется на \device~(именуемое в дальнейшем «устройство») и содержит сведения о его технических характеристиках, правилах эксплуатации, назначении и обслуживанию, а также необходимые сведения по функциональному назначению, основным параметрам, принципам работы, характеристикам, функциональным схемам формирования сигналов, перечню уставок и настраиваемых параметров.
\medskip
\par
До включения устройства в работу необходимо ознакомиться с настоящим руководством.
\medskip
\par
Надежность и долговечность устройства обеспечиваются не только его качеством, но и соблюдением режимов и условий эксплуатации, требований по транспортированию, хранению, монтажу. Поэтому выполнение всех требований, изложенных в настоящем документе, является обязательным. 
\medskip
\par
К эксплуатации устройства допускаются лица, изучившие настоящее руководство, прошедшие проверку знаний правил техники безопасности и эксплуатации электроустановок электрических станций и подстанций и имеющие группу по электробезопасности не ниже III по работе с электроустановками напряжением до 1000 В.

\vskip 75mm
\par
Компания Юнител Инжиниринг оставляет за собой авторские права на данный документ и на информацию, содержащуюся в нем, включая права на использование патентов. Копирование, использование и передача информации третьим лицам без письменного разрешения компании категорически запрещены.

\vskip 10mm
\par
По всем интересующим вопросам можно обращаться:
\begin{itemize}
\item по адресу электронной почты: \href{mailto:rza@uni-eng.ru}{\color{unimain}\uline{rza@uni-eng.ru}};
\item по телефону: +7 (495) 165-99-98 (доб. 2561).
\end{itemize}
Также с нашей продукцией можно ознакомиться на сайте: \href{https://uni-eng.ru}{\color{unimain}\uline{https://uni-eng.ru}}.